This project extends the baseline work of Sabir and Alebrahim (2025) on multi-threaded LU factorization by adding a comprehensive comparison with distributed parallel processing models. The research addresses the fundamental question: how do multi-threaded (scale-up) and distributed (scale-out) approaches compare for matrix workloads of increasing size on equal total core counts?

Using Python-based implementations with threading, multiprocessing, and Dask distributed frameworks, this study systematically measures completion time, peak memory consumption, and CPU utilization efficiency across matrix sizes from 200×200 to 2000×2000. The experimental design controls for total vCPU count, ensuring a fair comparison between parallelism models rather than hardware configurations.

Results from the committed \textbf{local} Dask \texttt{LocalCluster} suite (\texttt{summary\_statistics.json}; 90 configs × 3 iterations) demonstrate that multi-threaded approaches remain faster than local distributed execution for matrices from 200×200 through 2000×2000 at matched worker counts. No crossover appears in that local evidence set. Descriptive mean$\pm$std summaries are provided; inferential tests were planned but not coded in the artefact.

A matched-vCPU AWS EC2 round-1 was also collected (\texttt{results/live/ec2\_round1\_summary.json}): \textbf{1× t3.small} scale-up vs \textbf{2× t3.micro} scale-out (aggregate 2 vCPU). Scale-up numpy matmul at $n{=}500$ measured $\approx$0.0067\,s; on-node Dask \texttt{LocalCluster} per micro measured $\approx$0.52\,s mean. Multi-instance Dask registered two workers and passed a futures smoke test, but \texttt{da.matmul} \textbf{timed out}—no completed multi-instance matmul wall-clock is claimed.

Key contributions include: (1) direct performance comparison of threading vs. distributed models at equal core counts, (2) comprehensive metric collection (time, memory, CPU), (3) reproducible benchmarking framework with full test coverage, and (4) analysis methodology applicable to production cloud infrastructure decisions.
